\documentclass[11pt]{article}
\usepackage[a4paper,margin=28mm]{geometry}
\usepackage{amsmath,amsthm}
\usepackage[hidelinks]{hyperref}
\newtheorem{theorem}{Theorem}
\title{The Satisfiability Hypothesis from Clique by graph order}
\author{Atlas\thanks{The proof in this note was found using ChatGPT.}}
\date{18 November 2026}
\begin{document}
\maketitle

Let \textsc{Clique} ask whether a finite simple undirected graph with
individually represented vertices contains at least \(k\) pairwise
adjacent vertices. The threshold \(k\) is a nonnegative integer encoded
in binary. Write \(n\) for the number of vertices and \(B\) for the total
input length in bits. The Satisfiability Hypothesis (SH) excludes
deterministic CNF satisfiability algorithms in time
\(2^{O(s^b)}s^{O(1)}\) for each fixed \(0<b<1\), where \(s\) is the
normalized input length.

\hypertarget{clique-order-sh-converse}{}
\begin{theorem}
Suppose that for every fixed \(0<a<1\), \textsc{Clique} admits no
deterministic algorithm with running time \(2^{O(n^a)}B^{O(1)}\).
Then SH holds.
\end{theorem}

\begin{proof}
Suppose SH fails and fix a CNF decision algorithm of running time
\(2^{O(s^b)}s^{O(1)}\) for some fixed \(0<b<1\).
Quantitative sparsification for width three~\cite[Lemma~2.2]{sh}
provides a constant \(D\ge1\) such that, for each fixed \(0<d<1\),
a width-three formula on \(v\) variables is a disjunction, for
satisfiability purposes, of at most \(2^{O(v^{1-d})}\) width-three
formulas, each using at most \(v\) variables and
\(O(v^{1+Dd})\) clauses. This operation takes deterministic time
\(2^{O(v^{1-d})}\) times a fixed polynomial in the complete input
length. This quantitative lemma is a preprint dependency.

Relabeling variables consecutively gives every sparse output length
\(O((v+1)^{1+Dd}\log(v+2))\) bits. Choose \(d\) small enough that
\(b(1+Dd)<1\), and then choose \(a<1\) strictly larger than
\(\max\{1-d,b(1+Dd)\}\). Running the CNF solver on all outputs and taking
the disjunction gives one deterministic width-three satisfiability
algorithm with time \(2^{O(v^a)}\) times a fixed polynomial in the
complete formula length. The strict choice of \(a\) absorbs the
logarithmic factor in the encoding length.

For the graph reduction, dispose of \(k>n\) and \(k=0\) directly.
Use one Boolean variable \(x_u\) for every vertex \(u\). Add the clause
\(\neg x_u\lor\neg x_w\) for every distinct nonadjacent pair \(u,w\).
Use a Boolean circuit to force the sum of the selection bits to be
at least \(k\), as in the linear-variable approach of
Impagliazzo, Paturi and Zane~\cite[Theorem~5]{ipz}. More explicitly,
pad the input bits to the next power of two and add them in a balanced
tree using binary ripple-carry additions. Level \(j\) adds
\(O(n/2^j)\) pairs of \(O(j)\)-bit numbers. The circuit therefore uses
\(O(n\sum_{j\ge1}j/2^j)=O(n)\) constant-fanin gates, including comparing
the total with \(k\). Encode each fanin-two gate by one output variable
and a width-at-most-three clause for every disallowed input-output
triple. Unit clauses can fix constants and require the final comparison
output to be true. The resulting width-three formula has
\(O(n)\) variables, is constructible in time polynomial in \(B\), and is
satisfiable exactly when the selected vertices can form a clique of size
at least \(k\).

The width-three solver therefore decides \textsc{Clique} in time
\(2^{O(n^a)}B^{O(1)}\), contradicting the exclusion at exponent \(a\).
\end{proof}

\begin{thebibliography}{9}
\bibitem{sh}
MIT Hardness Group, J.~Brunner, E.~D. Demaine, F.~Frei, T.~Gomez,
M.~Hecher, J.~Li, and J.~Lynch.
\emph{ETH is Dead, Long Live SH}.
\href{https://arxiv.org/abs/2611.02650v2}{arXiv:2611.02650v2}, 2026.
\bibitem{ipz}
R.~Impagliazzo, R.~Paturi, and F.~Zane.
Which problems have strongly exponential complexity?
\emph{Journal of Computer and System Sciences}, 63(4):512--530, 2001.
\href{https://doi.org/10.1006/jcss.2001.1774}{DOI}.
\end{thebibliography}
\end{document}
